\chapter*{概要}
\thispagestyle{empty}
%
本研究では，軟弱地面における二足歩行ロボットの歩行安定性向上を目的として，重心高さが床変形の影響を受けて変動する環境に適応可能な歩行生成アルゴリズムを提案する．従来の歩行生成では，線形倒立振子モデル（LIPM）を用い，重心高さを一定と仮定することで重心運動を線形化し，ZMPやCapture Pointなどの指標に基づく制御設計を可能としてきた．この枠組みにより，平坦かつ剛体的な床面を仮定した環境においては，計算効率の高い歩行生成と安定な歩行制御が実現されてきた．一方，砂地やスポンジ床のような軟弱地面では，足底の沈み込みによって支持脚高さが変動し，重心高さ一定というLIPMの仮定が成立しなくなる．その結果，重心運動の予測誤差が増大し，従来のLIPMベース制御では歩行が不安定化しやすいという問題がある．

本研究では，床沈み込み量を推定し，その変動成分を重心高さの参照値に反映させるLIPM拡張型の重心高さ補償アルゴリズムを提案する．沈み込みが検出された場合には鉛直方向の速度および加速度項を目標重心高さに加算することで，動的に変化する床高さに追従可能な歩行軌道を生成する．この方法はLIPMの枠組みを維持しつつ鉛直方向の自由度を扱うため，計算コストが低く既存の歩行制御器へ容易に統合できる．

提案手法の有効性を検証するため，Choreonoid上に構築したCHIDORIロボットモデルを用いて，床をバネ・マス・ダンパ系で表現した飛び石地形シミュレーションを実施した．その結果，従来のLIPM制御では沈み込み発生時に重心が支持多角形外へ発散し転倒したのに対し，提案手法では鉛直方向合力の適切な増加と体幹の支持脚側への傾斜調整により重心発散が抑制され，歩行を継続できることを確認した．

以上より，本研究で提案する重心高さ補償アルゴリズムは，軟弱地面に固有の鉛直方向の不確実性に対してロバストに機能し，二足歩行ロボットの実環境適応性を高める手法として有効であることを示す．